\section{Results status}

\begin{center}
\fcolorbox{black}{gray!10}{\parbox{0.91\linewidth}{
\textbf{No protected scientific-superiority result is claimed.} The bounded
execution-integrity studies now support a narrower statement: cryptographic
receipts can fail closed against registered artifact corruption and can make
independent key domains load-bearing, but a valid artifact does not attest that
\emph{this} process invocation completed and a valid signature chain does not
establish scientific admission.}}
\end{center}

The results below are therefore organized by what the execution layer can and
cannot establish. They are measurements within frozen finite protocols, not a
general reliability estimate and not an authorization to collapse provenance,
execution integrity and scientific validity into one decision.

\subsection{Attestation composition is not scientific admission}
\label{sec:attestation-composition}

The frozen three-role attestation study verified Ed25519 chains under the
registered truncation, substitution, splice, reorder, cross-occurrence replay,
and stale-representation attacks. The central adverse control was full key-set
compromise: when the attacker possessed every role key, cryptographic validation
detected \textbf{0/6} forged cases and the unscoped chain-as-science reading
false-promoted all \textbf{6/6}. The experiment therefore establishes a boundary,
not a rescue: a signature proves possession of a signing key under the registered
verification rule; it does not prove that key custody was independent or that the
scientific conclusion was warranted.

The same study retained non-vacuity controls. Valid chains were accepted, the
registered scientific fields did not leak into the cryptographic/provenance layer,
and the scientific-admission decision remained separate from chain validity. The
current paper keeps that separation as a premise of every later experiment.

\subsection{Artifact corruption is detected without rejecting benign re-encoding}

A prospective execution-integrity successor tested six externally applied artifact
faults: truncation, signature-byte corruption, dropped chain link, public-key
substitution, favourable terminal relabelling and count inflation. All six faults
were applied and all six were rejected by the checker, with \textbf{0/6 false
promotions}. This result concerns committed-artifact integrity; it does not by
itself cover an error occurring inside the producer before the artifact is
committed.

The corresponding false-rejection control applied six semantics-preserving
re-encodings of a valid result. All six remained accepted, with \textbf{0/6 false
rejections}; four variants changed the serialized bytes. The checker therefore
canonicalizes the registered semantics rather than demanding raw byte identity.
This endpoint is deliberately narrow: it does not estimate false rejection under
clock skew, memory pressure, filesystem faults or other transient host conditions.

\subsection{Chain length is not the measured security axis}

The frozen trust-domain experiment separates two quantities that the earlier
three-link study confounded. At one trust domain, attack detection was
\textbf{1.000} for chain lengths $k=1,2,3$, with all four clean chains accepted at
each length. Increasing the number of signed links therefore produced no measured
improvement on the registered attack family. Longer chains may carry more semantic
roles, but length alone did not improve this endpoint.

At fixed chain length $k=3$, the experiment then varied the number of independent
cryptographic key domains while compromising exactly one domain. On four real
receipts, a one-domain construction accepted forged favourable chains in
\textbf{4/4} cases. With two independent domains the same attack was accepted in
\textbf{0/4}, and with three domains in \textbf{0/4}. Thus the measured false-promotion
rate stepped from $1.00$ to $0.00$ between $d=1$ and $d=2$ under this threat model,
whereas the $k$-only experiment was flat.

This is a cryptographic-independence result, not a theorem about organizational
independence. Independent key generation is the modeled property; separate teams,
institutions, machines and administrative domains are possible implementations but
were not all measured. The panel has four real receipts, one compromised domain and
$k=3$ for the $d$ comparison. No claim is made about multiple-domain compromise or
arbitrary trust organizations.

\subsection{Content integrity does not attest process liveness}

Host/process fault injection exposed a different boundary. When the producer was
terminated by \texttt{SIGKILL} or \texttt{SIGTERM} while an older valid artifact
remained at the output path, the artifact checker correctly returned green for that
old artifact. Two such stale-artifact cases were accepted. By contrast, an
unwritable output path and an absent output directory failed loudly and did not
silently produce a new result.

The implication is operationally important: an artifact verifier answers
``is this artifact valid?'', not ``did this invocation complete?''. A workflow that
uses a green artifact check as evidence of current-run success needs an additional
occurrence/liveness binding supplied by the orchestrator. This is not a checker bug;
it is the boundary of the object being checked.

\subsection{The measured cost is small in absolute time and large relative to a tiny baseline}

A no-attestation comparator reused the same frozen facts and canonicalization while
removing Ed25519 key derivation and signing. Across four cases and 40 repeats, median
composition time was approximately \textbf{0.105 ms} without attestation and
\textbf{1.444 ms} with attestation, an added median of \textbf{1.340 ms}. Because the
baseline is very small, the relative multiple is \textbf{13.8$\times$}; the measured
increment is about \textbf{111.6 $\mu$s per signed link}. Serialized output grew from
6926 to 9650 bytes, or \textbf{1.39$\times$}.

These figures measure composition-side cost only. Verification, I/O and orchestration
are excluded, so the numbers are not an end-to-end performance characterization.

\subsection{Provenance interoperability preserves the boundary}

A separate frozen study serialized execution-only facts through W3C PROV-JSON and an
RO-Crate workflow-run projection. Across the registered 22-case corpus, both
round-trips preserved the normalized execution vector, scientific fields remained
outside the provenance-only records, and imported versus native execution-status
interpretation agreed on every case. The representation therefore transports the
execution/science boundary without requiring an implementation-specific receipt
format.

The serialized sizes are descriptive rather than performance claims. They do not
support throughput, latency or storage projections at corpus scale.

\subsection{What the combined results establish}

The combined evidence supports four bounded conclusions. First, the registered
artifact checker rejects the tested committed-artifact corruptions without rejecting
the tested semantics-preserving encodings. Second, chain length and cryptographic
trust-domain count are distinct variables, and only the latter changed compromise
resistance in the registered panel. Third, artifact integrity and process liveness
are distinct properties. Fourth, none of these properties is scientific warrant:
a cryptographically and structurally valid receipt may still rest on a false premise,
a weak study design or a shared trust root.

Accordingly, any publication table must report the threat model, the number of trust
domains, the compromise pattern, false promotions, false rejections, liveness status
and the scientific-admission disposition as separate fields. A future general
reliability claim would require a larger independently custodied campaign rather than
pooling these bounded protocol checks into one rate.
